\section*{Scope, reading path, and headline result}

This paper is about using numerical statistical mechanics to audit collective
behavior \emph{inside learning systems}. It is not a survey of machine learning
for discovering phases of physical matter, and it does not claim that every
learning curve has a thermodynamic limit. Readers interested in the conceptual
contribution can follow Sections~\ref{sec:frontier}--\ref{sec:questions}; readers
interested in estimators and finite-size logic can continue through
Sections~\ref{sec:method}--\ref{sec:statistics}; and readers interested in the
measurements can start at Sections~\ref{sec:design}--\ref{sec:results}.

The headline empirical result is deliberately negative. Increasing width and
data improves a small Transformer, and matched MLP ablation has a measurable
effect, but the reference grid fails the registered tests needed for phase
language. The finding is \emph{learning response observed; phase boundary not
established}. This bounded conclusion is the intended behavior of the atlas,
not a failure to produce a positive plot.

\manuscriptfigure{figures/figure1_atlas_taxonomy.pdf}{The atlas has three
layers. Seven scientific coordinates specify what object is under study. A
protocol specifies how the object is measured, perturbed, and held out.
Independent theory, fidelity, support, and outcome fields specify what may be
claimed. Collapsing these layers is the main source of ambiguous phase language
in machine-learning experiments.}{fig:atlas-overview}

\FloatBarrier
